% =============================================================================
% AVALIAÇÃO QUALITATIVA E MÉTODOS MISTOS (16:55 a 17:50)   [Exercício 3]
% =============================================================================
\section[Qualitativo]{Avaliação qualitativa e métodos mistos}

\begin{frame}{O que os números não mostram}
  O DiD do \textbf{Qualifica} mostrou \destaque{+R\$ 300} em média. Mas:
  \vspace{0.2cm}
  \begin{enumerate}
    \item O efeito é \textbf{muito menor} entre mulheres com filhos pequenos (R\$ 80);
    \item \textbf{35\%} dos inscritos abandonaram o curso;
    \item \textbf{80\%} se dizem ``muito satisfeitos'', mas o gestor acha que ``algo não bate''.
  \end{enumerate}
  \vspace{0.3cm}
  \begin{block}{As perguntas que sobram}
    \small \emph{Por quê?} \emph{Para quem?} \emph{O que significa ``satisfeito''?} Os números identificam o efeito; \ind{a avaliação qualitativa ajuda a explicá-lo}.
  \end{block}
\end{frame}

\begin{frame}{Quantitativo $\times$ qualitativo}
  \begin{columns}[T]
    \begin{column}{0.5\textwidth}
      \begin{block}{Quantitativo}
        \footnotesize
        \begin{itemize}
          \item \textbf{Quanto?}
          \item Efeito \emph{médio}
          \item Contrafactual, viés de seleção
          \item Generalização e replicação
        \end{itemize}
      \end{block}
    \end{column}
    \begin{column}{0.5\textwidth}
      \begin{exerciciobox}{Qualitativo}
        \footnotesize
        \begin{itemize}
          \item \textbf{Como e por quê?}
          \item O \emph{específico}: contexto e significados
          \item Poucos casos, em profundidade
          \item Percepções, valores, mecanismos
        \end{itemize}
      \end{exerciciobox}
    \end{column}
  \end{columns}
  \vspace{0.3cm}
  \centering\small Não é ``quali \emph{ou} quanti'': são \textbf{perguntas diferentes}.
\end{frame}

\begin{frame}{Para que serve a avaliação qualitativa?}
  \begin{enumerate}
    \item \textbf{Reconstruir a política} a partir dos atores: como ela foi \emph{de fato} implementada e percebida;
    \item \textbf{Identificar impactos diferenciados}: quem ganha, quem não ganha, e por quê;
    \item \textbf{Captar o que é difícil de medir}: empoderamento, dignidade, confiança, capital social;
    \item \textbf{Explicar mecanismos}: por que a política funciona em alguns contextos e em outros não (``abrir a caixa-preta'').
  \end{enumerate}
\end{frame}

\begin{frame}{O qualitativo em cada etapa da avaliação}
  \centering
  \begin{tikzpicture}[scale=0.86,every node/.style={transform shape}]
    \node[cxi,text width=2.2cm] (a) at (0,0) {\textbf{Diagnóstico}\\{\scriptsize ouvir o público-alvo; entender o problema}};
    \node[cxi,text width=2.2cm] (b) at (3.1,0) {\textbf{Implementação}\\{\scriptsize observação; gestores e burocratas de rua}};
    \node[cxl,text width=2.2cm] (c) at (6.2,0) {\textbf{Resultados e impacto}\\{\scriptsize percepções, mecanismos, efeitos inesperados}};
    \node[cxi,text width=2.2cm] (d) at (9.3,0) {\textbf{Uso dos achados}\\{\scriptsize devolutiva e deliberação}};
    \draw[seta] (a) -- (b); \draw[seta] (b) -- (c); \draw[seta] (c) -- (d);
  \end{tikzpicture}
  \vspace{0.3cm}\par
  \small Hoje o foco é a etapa de \textbf{impacto}, mas o qualitativo também aparece antes e depois.
\end{frame}

\begin{frame}{Desenhando uma pesquisa qualitativa}
  \centering
  \begin{tikzpicture}[scale=0.82,every node/.style={transform shape}]
    \node[cxi,text width=2.0cm] (a) at (0,0) {\textbf{1. Pergunta}};
    \node[cxi,text width=2.0cm] (b) at (2.6,0) {\textbf{2. Casos e amostra}};
    \node[cxi,text width=2.0cm] (c) at (5.2,0) {\textbf{3. Instrumentos}};
    \node[cxi,text width=2.0cm] (d) at (7.8,0) {\textbf{4. Campo}};
    \node[cxl,text width=2.0cm] (e) at (10.4,0) {\textbf{5. Análise}};
    \draw[seta] (a) -- (b); \draw[seta] (b) -- (c); \draw[seta] (c) -- (d); \draw[seta] (d) -- (e);
    \node[font=\tiny,align=center,text width=2cm] at (0,-1.4) {o que queremos entender?};
    \node[font=\tiny,align=center,text width=2.2cm] at (2.6,-1.4) {amostra \emph{intencional}: variabilidade de posições};
    \node[font=\tiny,align=center,text width=2.2cm] at (5.2,-1.4) {roteiro, protocolo de observação};
    \node[font=\tiny,align=center,text width=2.2cm] at (7.8,-1.4) {tempo, confiança, ética (consentimento, anonimato)};
    \node[font=\tiny,align=center,text width=2.2cm] at (10.4,-1.4) {codificação, temas, validação};
  \end{tikzpicture}
  \vspace{0.3cm}
  \begin{itemize}
    \item A amostra qualitativa busca \ind{variabilidade} (quem teria uma visão diferente?), não probabilidade estatística;
    \item Pode-se escolher os casos com \textbf{evidência quantitativa} (ex.: maiores e menores efeitos).
  \end{itemize}
\end{frame}

\begin{frame}{Técnicas: visão geral}
  \scriptsize
  \renewcommand{\arraystretch}{1.3}
  \begin{tabularx}{\textwidth}{@{}>{\raggedright\arraybackslash}p{2.3cm}>{\raggedright\arraybackslash}X>{\raggedright\arraybackslash}X@{}}
    \toprule
    \textbf{Técnica} & \textbf{Serve para} & \textbf{Cuidados} \\ \midrule
    Observação, etnografia & Ver a política em ação, o cotidiano, o que não é dito & Tempo de campo; reflexividade \\
    Entrevista em profundidade & Experiências, significados, temas sensíveis & Roteiro; confiança; viés do entrevistador \\
    Grupo focal & Debate, consensos e divergências entre atores & Moderação; participantes dominantes \\
    Análise documental & Regras, metas, registros de implementação & Distância entre o papel e a prática \\
    Estudo de caso, \emph{process tracing} & Reconstruir sequências e testar mecanismos & Escolha dos casos \\ \bottomrule
  \end{tabularx}
\end{frame}

\begin{frame}{Observação e etnografia}
  \begin{itemize}
    \item \textbf{Imersão} na rotina do grupo; notas de campo e descrições densas;
    \item \textbf{Participante} ou \textbf{não participante}; quanto mais tempo, menor a \emph{reflexividade} (o grupo se comporta como sempre);
    \item Útil para ver \ind{como as pessoas reagem aos instrumentos} da política.
  \end{itemize}
  \vspace{0.2cm}
  \begin{block}{Adato (2008): transferências de renda na Nicarágua e na Turquia}
    \scriptsize
    O survey indicava um programa bem \emph{focalizado}, mas a comunidade via o critério como \textbf{arbitrário}. Mães não davam o suplemento de ferro por julgarem que fazia mal; crianças eram superalimentadas antes da pesagem obrigatória para ``passar no teste''.
  \end{block}
\end{frame}

\begin{frame}{Entrevistas em profundidade}
  \begin{columns}[T]
    \begin{column}{0.52\textwidth}
      \footnotesize
      \begin{itemize}
        \item A \emph{palavra} do participante é o foco; perguntas \textbf{abertas e flexíveis};
        \item \textbf{Estruturada}, \textbf{semiestruturada} (tópicos-guia), \textbf{não estruturada} ou \textbf{narrativa};
        \item Quantas? Até a \ind{saturação}: quando novas entrevistas deixam de trazer ideias novas.
      \end{itemize}
    \end{column}
    \begin{column}{0.48\textwidth}
      \begin{exerciciobox}{Exemplo: roteiro (Qualifica)}
        \scriptsize
        \begin{enumerate}
          \item Como você soube do curso e por que se inscreveu?
          \item Conte como era uma semana típica de curso.
          \item O que dificultou ou ajudou a continuar?
          \item O que mudou na sua vida depois?
        \end{enumerate}
      \end{exerciciobox}
    \end{column}
  \end{columns}
  \vspace{0.2cm}
  \footnotesize \alert{Evite} perguntas fechadas ou indutivas (``O curso foi bom, não foi?'').
\end{frame}

\begin{frame}{Grupos focais}
  \begin{columns}[T]
    \begin{column}{0.55\textwidth}
      \small
      \begin{itemize}
        \item \textbf{Não} é a soma de entrevistas: o foco está na \ind{discussão} entre os participantes;
        \item O moderador estimula quem fala pouco, contém quem domina e mantém o tema;
        \item Mostra \textbf{consensos e divergências}, normas e linguagem do grupo.
      \end{itemize}
    \end{column}
    \begin{column}{0.45\textwidth}
      \begin{block}{Hunter e Sugiyama (2014)}
        \scriptsize
        11 grupos focais em 3 municípios do Nordeste + entrevistas com gestores. O Bolsa Família aparece associado a \textbf{dignidade}, \textbf{empoderamento} e menor dependência do clientelismo local.
      \end{block}
    \end{column}
  \end{columns}
\end{frame}

\begin{frame}{Documentos, estudos de caso e \textit{process tracing}}
  \begin{itemize}
    \item \textbf{Análise documental}: leis, portarias, atas, relatórios, bases administrativas. Mostra o \emph{desenho formal}; a observação e as entrevistas mostram a \emph{prática};
    \item \textbf{Estudo de caso}: um município, uma escola, um serviço, analisados em profundidade;
    \item \textbf{\emph{Process tracing}}: reconstruir a sequência de eventos e testar, com evidências do próprio caso, se o \ind{mecanismo} proposto operou (Beach e Pedersen, 2013).
  \end{itemize}
  \vspace{0.3cm}
  \begin{block}{Combinação típica}
    \small Documentos (o que deveria acontecer) + observação (o que acontece) + entrevistas (o que significa para os atores).
  \end{block}
\end{frame}

\begin{frame}{Da fala ao tema: análise dos dados}
  \centering
  \begin{tikzpicture}[scale=0.85,every node/.style={transform shape}]
    \node[cx,text width=8.6cm,align=left] (t) at (0,1.9) {\textbf{Trecho:} ``O curso era à noite e eu não tinha com quem deixar as crianças.''};
    \node[cxi,text width=8.6cm,align=left] (c) at (0,0.7) {\textbf{Códigos:} falta de creche; horário noturno};
    \node[cxi,text width=8.6cm,align=left] (g) at (0,-0.5) {\textbf{Categoria:} barreiras de cuidado};
    \node[cxl,text width=8.6cm,align=left] (h) at (0,-1.7) {\textbf{Tema:} custos de acesso e de cuidado que afastam as mães do curso};
    \draw[seta] (t) -- (c); \draw[seta] (c) -- (g); \draw[seta] (g) -- (h);
  \end{tikzpicture}
  \vspace{0.2cm}\par
  \footnotesize Passos da \ind{análise temática}: transcrever e ler, codificar, agrupar, rever os temas, nomeá-los e relatar com trechos (Braun e Clarke, 2006).
\end{frame}

\begin{frame}{Rigor na pesquisa qualitativa}
  \begin{columns}[T]
    \begin{column}{0.5\textwidth}
      \begin{itemize}
        \item \textbf{Triangulação}: várias fontes e técnicas;
        \item \textbf{Transparência}: livro de códigos e trilha de decisões;
        \item \textbf{Reflexividade}: explicitar a posição do pesquisador.
      \end{itemize}
    \end{column}
    \begin{column}{0.5\textwidth}
      \begin{itemize}
        \item \textbf{Devolutiva} aos participantes;
        \item \textbf{Saturação} e busca de casos negativos;
        \item \textbf{Limites}: os achados não dão magnitude nem generalização.
      \end{itemize}
    \end{column}
  \end{columns}
  \vspace{0.3cm}
  \begin{block}{Críticas comuns}
    \small Subjetividade, ausência de regularidades, menor precisão e dificuldade de replicação: respondem-se com \emph{rigor explícito}, não com a imitação do quantitativo.
  \end{block}
\end{frame}

\begin{frame}{Métodos mistos: por que combinar?}
  \begin{enumerate}
    \item \textbf{Credibilidade}: convergência de achados (triangulação);
    \item \textbf{Compreensão ampla}: diferentes dimensões do mesmo fenômeno;
    \item \textbf{Descoberta} de fatores inesperados;
    \item \textbf{Diversidade} de perspectivas, inclusive discordantes.
  \end{enumerate}
  \vspace{0.3cm}
  \begin{block}{Regra de ouro}
    \small O método é escolhido pela \textbf{pergunta}, e não o contrário.
  \end{block}
\end{frame}

\begin{frame}{Desenhos de métodos mistos}
  \centering
  \begin{tikzpicture}[scale=0.82,every node/.style={transform shape}]
    \node[font=\footnotesize\bfseries,anchor=east] at (2.2,2.2) {Convergente};
    \node[cxi,minimum width=2.4cm] (a1) at (3.8,2.5) {QUANT};
    \node[cxl,minimum width=2.4cm] (a2) at (3.8,1.9) {QUAL};
    \node[cx,minimum width=3.6cm] (a3) at (7.8,2.2) {comparar e integrar};
    \draw[seta] (a1.east) -- (a3.west); \draw[seta] (a2.east) -- (a3.west);
    \node[font=\footnotesize\bfseries,anchor=east] at (2.2,0) {Sequencial explanatório};
    \node[cxi,minimum width=2.4cm] (b1) at (3.8,0) {QUANT};
    \node[cxl,minimum width=2.4cm] (b2) at (7.8,0) {QUAL};
    \draw[seta] (b1.east) -- node[above,font=\tiny]{explicar} (b2.west);
    \node[font=\footnotesize\bfseries,anchor=east] at (2.2,-1.8) {Sequencial exploratório};
    \node[cxl,minimum width=2.4cm] (c1) at (3.8,-1.8) {QUAL};
    \node[cxi,minimum width=2.4cm] (c2) at (7.8,-1.8) {QUANT};
    \draw[seta] (c1.east) -- node[above,font=\tiny]{construir} (c2.west);
  \end{tikzpicture}
  \vspace{0.2cm}\par
  \footnotesize Explanatório: o survey acha o padrão e as entrevistas o explicam. Exploratório: entrevistas orientam o questionário. (Creswell e Plano Clark)
\end{frame}

\begin{frame}{Métodos mistos na prática}
  \footnotesize
  \begin{itemize}
    \item \textbf{Alfabetização de adultos} (Patton, 2002). Survey: \textbf{77\% ``muito satisfeitos''}. As entrevistas mostraram por quê: ler os remédios dos filhos, a Bíblia, as compras. O ganho era \ind{empoderamento cotidiano}, não emprego.
    \item \textbf{Focalização no Chile} (Clert e Wodon, 2001). Survey: grande sucesso na focalização. Grupos focais: beneficiários não entendiam as regras e sentiam \alert{injustiça} na distribuição.
    \item \textbf{Desenvolvimento comunitário na Índia} (Bamberger, 2012). 200 comunidades sorteadas $\to$ pesquisa exploratória $\to$ survey $\to$ estudo em profundidade de 5 projetos tratados e 5 controles $\to$ triangulação.
  \end{itemize}
\end{frame}

\begin{frame}{Exercício 3: do trecho ao tema (20 min, em grupos)}
  \begin{exerciciobox}{\texttt{atividades/exercicio\_03\_codificando\_entrevistas.md}}
    \small
    \begin{enumerate}
      \item \textbf{Codifique} os 8 trechos (1 a 3 códigos cada);
      \item \textbf{Agrupe} em categorias e proponha \textbf{2 a 3 temas};
      \item Que quebra-cabeça do Qualifica cada tema explica? Quem mais entrevistar? Que outra técnica usar?
    \end{enumerate}
  \end{exerciciobox}
  \vspace{0.2cm}
  \small Salas simultâneas (20 min) e plenária (10 min). Os trechos estão nos próximos slides e no arquivo.
\end{frame}

\begin{frame}{Trechos (1 de 2)}
  \scriptsize
  \renewcommand{\arraystretch}{1.3}
  \begin{tabularx}{\textwidth}{@{}c>{\raggedright\arraybackslash}p{2.4cm}>{\raggedright\arraybackslash}X@{}}
    \toprule
    \textbf{\#} & \textbf{Quem fala} & \textbf{Trecho} \\ \midrule
    1 & Maria, 34, mãe de dois, desistiu & ``O curso era à noite e eu não tinha com quem deixar as crianças. Tentei levar a menor uma vez, mas a professora disse que não podia.'' \\
    2 & João, 27, concluiu e foi contratado & ``O certificado ajudou, mas o que valeu foi o professor ter me apresentado numa empresa onde trabalhava um ex-aluno.'' \\
    3 & Dona Lúcia, 52, concluiu, sem emprego & ``Aprendi muita coisa, mas as vagas pedem experiência e quando veem minha idade a conversa muda.'' \\
    4 & Carla, 31, desistiu & ``A passagem de ônibus até o curso custava quase o que eu ganhava num dia. Chegou uma hora que não compensava.'' \\ \bottomrule
  \end{tabularx}
\end{frame}

\begin{frame}{Trechos (2 de 2)}
  \scriptsize
  \renewcommand{\arraystretch}{1.3}
  \begin{tabularx}{\textwidth}{@{}c>{\raggedright\arraybackslash}p{2.4cm}>{\raggedright\arraybackslash}X@{}}
    \toprule
    \textbf{\#} & \textbf{Quem fala} & \textbf{Trecho} \\ \midrule
    5 & Rafael, 22, concluiu & ``Gostei da turma. A gente trocava contato e depois um indicava o outro pra algum bico.'' \\
    6 & Gestora municipal & ``Matriculamos 400 pessoas porque a meta era essa. Acompanhar quem saiu, não deu: não tínhamos equipe.'' \\
    7 & Instrutor & ``Metade da turma não tinha o básico de informática. Eu tinha que escolher: ensinar o conteúdo ou ensinar a ligar o computador.'' \\
    8 & Ana, 29, mãe, concluiu & ``Eu me sentia mais segura pra procurar emprego, mas a vaga que apareceu era em outro turno e a creche fecha às 17h.'' \\ \bottomrule
  \end{tabularx}\par
  \vspace{0.2cm}\tiny Trechos fictícios, inspirados em situações comuns em programas de qualificação.
\end{frame}
